% TOP_PROBLEM: {"problemId": "problem.polynomial-time-quasioptimal-logarithmic-energy-points-on-the-2-sphere", "canonicalTitle": "Polynomial-Time Quasioptimal Logarithmic-Energy Points on the 2-Sphere", "releaseRank": 474, "primaryDomain": "applied_computational_mathematics", "primaryDomainLabel": "Applied & computational mathematics", "displayStatus": "open", "releaseStatus": "open"}
% !TeX program = lualatex
% Definition and sources only; this file is not an LLM solution attempt.
\documentclass[11pt]{article}
\usepackage[margin=25mm]{geometry}
\usepackage{fontspec}
\setmainfont{Latin Modern Roman}
\usepackage{amsmath,amssymb,mathtools}
\usepackage{unicode-math}
\setmathfont{Latin Modern Math}
\usepackage{xurl,hyperref}
\hypersetup{colorlinks=true,urlcolor=blue}
\setcounter{secnumdepth}{0}
\setlength{\parindent}{0pt}
\setlength{\parskip}{5pt}
\setlength{\emergencystretch}{3em}
\begin{document}
\section*{\texorpdfstring{Polynomial-Time Quasioptimal Logarithmic-Energy Points on the 2-Sphere}{Polynomial-Time Quasioptimal Logarithmic-Energy Points on the 2-Sphere}}\label{474}
\subsection{Definitions and mathematical statement}
For distinct points $x_1,\ldots,x_N$ on the unit sphere $S^2\subset{\mathbb{R}}^3$, define logarithmic energy, using Smale's pair-counting convention, by
\[
 E(x_1,\ldots,x_N)=\sum_{1\le i<j\le N}
 \log\frac1{\|x_i-x_j\|},\qquad E_N^{\min}=\inf E.
\]
Seek a uniform real-number algorithm and an absolute constant $c$ such that, for every $N\ge2$, it produces $N$ points in time polynomial in $N$ and
\[
 E(x_1,\ldots,x_N)\le E_N^{\min}+c\log N.
\]
The real-number computation model and admissible operations are those in Smale's Problem 7. This does not silently impose a bit-complexity bound for exact transcendental coordinates. The required error is additive and logarithmic, much smaller than a fixed relative error in the leading-order energy.
\par\smallskip\textit{Formulation and concept references:} \href{https://www.cityu.edu.hk/ma/doc/people/smales/pap104.pdf}{[S1]}.

\subsection{Short English statement}
Can a polynomial-time real-number algorithm place points on the sphere with logarithmic energy within a constant times log N of optimal?
\subsection{Sources}
\begin{itemize}
\item[S1]Stephen Smale, Mathematical Problems for the Next Century, Problem 7 (1998).
\url{https://www.cityu.edu.hk/ma/doc/people/smales/pap104.pdf}.
\textit{Formulation source.}
\end{itemize}
\end{document}
